\section{Arquitetura e Parametrização do Modelo Simscape Specialized Power Systems}\label{sec:modelo_simscape}

A investigação dos transitórios eletromecânicos no domínio do tempo exige modelos matemáticos não lineares formulados no referencial ortogonal $dq0$, capazes de capturar a dinâmica de acoplamento dos fluxos magnéticos e as oscilações de conjugado sob perturbações na rede ou no eixo. A correta parametrização do bloco \textit{Asynchronous Machine SI Units} no ambiente Simscape SPS depende da aplicação rigorosa do invariante de transformação de impedâncias Estrela-Triângulo para preservar a fidelidade física das grandezas terminais.

% =========================================================================
% [DIRETRIZ DE REDAÇÃO]
% - Objetivo: Atender ao Item 2.A de sim_proj.pdf. Apresentar o modelo dinâmico implementado
%   no Simscape Specialized Power Systems (SPS) utilizando o bloco Asynchronous Machine SI Units.
% - Formulação: Apresentar a dedução formal do invariante de mapeamento Estrela-Triângulo (ADR-002):
%   R_{1,Y} = R_1 / 3, R_{2,Y}' = R_2' / 3, L_{l1,Y} = L_{l1} / 3, L_{l2,Y}' = L_{l2}' / 3, L_{m,Y} = L_m / 3.
% - Figuras/Tabelas: Inserir a tabela de parâmetros mapeados e o diagrama esquemático do modelo.
% - Confronto: Justificar a configuração do solver ode23tb (TR-BDF2, L-estável) para assegurar
%   estabilidade numérica sob correntes de inrush e degraus abruptos de conjugado mecânico.
% =========================================================================

\subsection{O Invariante de Mapeamento Estrela-Triângulo ($\text{Y}-\Delta$)}\label{subsec:mapeamento_ydelta}

O bloco \textit{Asynchronous Machine SI Units} do Simscape adota internamente a representação por enrolamentos em Estrela ($\text{Y}$) não aterrada, recebendo como entrada a tensão eficaz entre linhas da rede ($V_L = \SI{220}{\volt}$). Para modelar com exatidão a máquina ensaiada conectada fisicamente em Triângulo ($\Delta$) sem alterar as tensões nodais dos barramentos trifásicos, aplica-se a transformação de impedâncias:
\begin{equation}
R_{1,Y} = \frac{R_{1,\Delta}}{3} = \frac{\SI{1,189}{\ohm}}{3} = \SI{0,39633}{\ohm}, \qquad R_{2,Y}' = \frac{R_{2,\Delta}'}{3} = \frac{\SI{1,180}{\ohm}}{3} = \SI{0,39333}{\ohm},
\label{eq:resistencias_estrela}
\end{equation}
\begin{equation}
L_{l1,Y} = \frac{L_{l1,\Delta}}{3} = \SI{1,90833}{\milli\henry}, \quad L_{l2,Y}' = \frac{L_{l2,\Delta}'}{3} = \SI{2,80667}{\milli\henry}, \quad L_{m,Y} = \frac{L_{m,\Delta}}{3} = \SI{43,85700}{\milli\henry}.
\label{eq:indutancias_estrela}
\end{equation}

Os parâmetros mecânicos acoplados ($J = \SI{0,701}{\kilogram\meter\squared}$, $B = \SI{0,0426}{\Nm\cdot\second/\radian}$ e $p_{\text{pares}} = 2$) permanecem invariantes. A equivalência física garante que as correntes de linha, o fluxo de potência ativa e reativa e o conjugado eletromagnético gerados no Simscape coincidam rigorosamente com a operação em Triângulo.

\subsection{Configuração do Solver Numérico L-Estável}\label{subsec:config_solver}

Devido à elevada rigidez numérica (\textit{stiffness}) decorrente das constantes de tempo elétricas subtransitórias frente às constantes de tempo mecânicas da inércia pesada, adotou-se o solver numérico L-estável \texttt{ode23tb} (combinação trapezoidal e diferenciação retroativa TR-BDF2) com tolerância relativa de $10^{-4}$ e passo máximo de amostragem de $\SI{1e-4}{\second}$.
